\section{Discussion}

Across three cohorts of severe mental illness, a single Bayesian measurement pipeline
places clinical and biological measurements in one transdiagnostic coordinate system,
holds that eight-dimension map fixed, and validates it for structure, temporal coherence
and prognosis. The contribution is not a new estimator but the discipline of its
integration, novel in combination on three counts. The measurement layer is
\emph{missingness-native}: one global observed-cell likelihood carries the modelled pool's
$\approx 40\%$ mean cell-missingness and two-cohort indicators, fit jointly rather than
imputed or deleted. \emph{Uncertainty is propagated end-to-end}: each patient's posterior
coordinates and their covariance flow into structure-testing, temporal decomposition and
errors-in-variables prognosis rather than collapsing to point scores at the measurement
boundary. And the map is \emph{fixed with diagnosis held out, then validated}---frozen
before any validation stage touches it, which is what licenses the downstream results as
properties of the measurement rather than artifacts fitted to an outcome. The
immunometabolic dissociation that runs through the paper is the existence proof for that
discipline: the same rigor that makes the map honest reveals structure a looser
measurement would blur---one recurring empirical property in which immunometabolic load is
the domain least entangled with overall functional severity. The cardiometabolic and inflammatory markers resolve as a single
immunometabolic axis anchored by body composition (BMI loading $\approx 0.95$, weight
$0.88$, waist circumference $0.84$), with smaller metabolic, inflammatory and
cardiovascular contributions (C-reactive protein $\approx 0.37$); that axis sits nearly
independent of \Gfac{}. Because the one inflammatory marker available here, CRP, is itself
substantially adiposity-driven~\citep{kappelmann2021,norton2025crp}, we read this as an
adiposity-anchored cardiometabolic axis and do not interpret it as an inflammatory
dimension independent of body composition---the IL-6-type signal implicated in mood
pathology~\citep{ye2021inflammation,milaneschi2021} was not measured here. The axis is
nonetheless more than body habitus: excluding the anthropometric indicators and refitting,
a coherent cardiometabolic--inflammatory axis survived, highly congruent with the full axis
(Tucker $\varphi=0.90$; Supplement~\ref{annA:nobmi}). If biological burden simply tracked how impaired a
patient is, it would add little to a clinician's global impression. The data say the
opposite---biological risk rides on an axis the severity picture does not see---and the
five archetypal extremes make the consequence concrete. The immunometabolic corner and
the severe, clean-biology corner carry comparable functional burden but opposite
biological profiles, and they diverge in prognosis (the immunometabolic corner is the
worst-prognosis pole and the well pole the best, two-year functional remission spanning
$22\%$ to $63\%$ across the five corners). This is the precise property a
biology-aware stratification can exploit: it separates patients who look equally ill
but are biologically opposite, and that separation is durable and prognostically
meaningful.

This headline is corroborated, in direction, by a large and independent literature,
which is why we state it as relative rather than absolute independence. Metabolic
burden in severe mental illness is driven by medication, adiposity, age and chronicity
more than by diagnosis or symptom severity, is present before chronic severity
accrues, and is comparable across schizophrenia, bipolar disorder and
depression~\citep{vancampfort2015mets,vancampfort2016diabetes,pillinger2017,perry2019,garridotorres2022};
inflammation marks a subgroup and specific symptoms rather than the severity
continuum~\citep{osimo2019,fernandes2016bd,fernandes2016sz}; and the
immuno-metabolic-depression programme has independently arrived at a separable
biological dimension mapping onto atypical, energy-related
symptoms~\citep{penninx2024imd,milaneschi2020,lamers2013,lamers2020,milaneschi2017}.
That a completely different method, in different cohorts, converges on a separable
biology axis is strong convergent validity. We also state the boundary conditions
plainly: case--control inflammatory elevations are
real~\citep{howren2009,osimo2020,haapakoski2015}, a modest severity or duration
dose--response for metabolic burden exists~\citep{penninxlange2018,pillinger2022},
metabolic dysfunction does track cognition and functioning in mood
disorders~\citep{maksyutynska2024}, and much of the inflammation--symptom association
is carried by adiposity rather than inflammation \emph{per
se}~\citep{kappelmann2021,zagaria2024}. Our small immunometabolic--\Gfac{} couplings
are best read as a quantification of a known qualitative pattern, sitting at its low
end because \Gfac{} indexes \emph{functional burden}---from which biology is most
decoupled---rather than the case--control contrast of being ill. That they survive,
and even fall slightly under, adjustment for medication, adiposity and site is the
strongest available evidence that the dissociation is not a confound.

What is new is less the direction of this finding---which the literature above
corroborates---than where biology sits in the model. The dimensional-psychopathology
tradition builds its coordinate systems from symptom report alone, with no biological
measurement inside the space~\citep{kotov2017hitop}, and normative modelling charts
per-measure deviation from a reference rather than a shared low-dimensional
map~\citep{marquand2016,wolfers2018}. The closest biological precedent identifies a
separable immuno-metabolic profile \emph{within} major depression, anchored on symptom
dimensions with biology entered as a correlated \emph{outcome}~\citep{lamers2013,lamers2020}, while
the B-SNIP biotype programme resolves discrete brain-based groups rather than a
continuum~\citep{clementz2016biotypes}. Placing routine biology as a co-equal latent axis
\emph{inside} a single diagnosis-blind map across bipolar disorder, schizophrenia and
depression---and testing its decoupling, durability and prognosis in that same
model---is the slot none of these occupy.

A second recurring result is structural. On these clinical--biological coordinates the
patient space is a graded continuum, not a set of discrete biotypes, and it cuts
almost entirely across DSM-5. This aligns with normative-modelling work reporting
extensive individual heterogeneity in the same
disorders~\citep{wolfers2018,marquand2016,segal2023}, and it stands in apparent
tension with the brain-based B-SNIP biotypes~\citep{clementz2016biotypes,clementz2020}.
The tension is more apparent than real: the two programmes interrogate different
measurement spaces, B-SNIP itself found DSM diagnoses to be a severity continuum, and
recent syntheses endorse both categorical and dimensional
views~\citep{clementz2024}; clustering choices can also resurface apparent biotypes
from continuous data~\citep{pan2026}. We therefore claim no biotypes \emph{in this
eight-dimensional clinical--biological space}---a claim we make rigorous with a
single-Gaussian falsification null (best-partition silhouette $0.140$ against a
structureless-Gaussian null of $0.137 \pm 0.002$, $z = 1.13$) rather than by the
absence of an elbow---not that biotypes are absent from all spaces. The five
archetypes are corners of this continuum, not clusters within it; they cut almost
orthogonally across DSM-5 (adjusted Rand index $0.021$) yet describe a patient's position
on the eight-dimensional map about $9.7\times$ more tightly than diagnosis does (mean
$\eta^2$ $0.256$ versus $0.026$). The corners span a severity gradient \emph{and} a
biological one---the immunometabolic and severe, clean-biology corners are equally severe
but biologically opposite---so the simplex separates on profile as well as severity; the
sharper claim that a partition is driven by activation and suicidality rather than overall
severity belongs to the coarse two-region tessellation (Supplement~\ref{ann:stratification}),
a distinct, lower-resolution object. Anchoring \Gfac{} on functioning alone also
sidesteps the well-known instability in what a symptom-only general factor
means~\citep{heinrich2021bifactor,levinaspenson2020,watts2021}, while remaining
consistent with the impairment reading of the \textit{p}
factor~\citep{caspi2014pfactor,lahey2017general}.

The remaining contributions are deliberately calibrated. The map persists over two
years---durable biology, moving symptoms---which both validates the measurement
longitudinally and motivates a clinical division of labour between traits to stratify
on and states to monitor. The immunometabolic axis is the single most durable
dimension (test--retest ICC $0.91$, fully invariant across visits; it remains trait-like,
ICC $0.76$, even with the anthropometric indicators removed, so the durability is not merely
body mass's stability---Supplement~\ref{annA:nobmi}), cognition is also
trait-like, while severity and symptoms slide as the population improves: the durable
part of the map is precisely its biological part.

This division is itself a design principle---\emph{stratify on what stays put; monitor
what moves}. Because the immunometabolic axis is both the durable trait and the
dimension least entangled with functional burden, it is the stable axis on which to fix a
patient's position---for instance at the immunometabolic corner (\textbf{A2}); general
burden is durable only \emph{by rank} (ICC $0.62$), with patients largely keeping their
relative position even as the cohort improves in absolute terms. The symptom axes, by
contrast, slide---population shifts of $-0.46$ on severity and $-0.84$ on suicidality from
baseline to two years, against an essentially static $+0.04$ on biology---and are
therefore quantities to track over time rather than to fix a stratum on. We frame this as
an interpretive design principle, read off the baseline structure and its forward-scored
longitudinal coherence rather than from any validated protocol: what the durability buys
is a principled separation of the axes worth stratifying on from those worth monitoring,
not an individual treatment rule.

The map predicts future functioning
incrementally beyond diagnosis and severity (archetypes $\Delta$ELPD $+62.8$
held-out, robust to inverse-probability weighting and permutation), but the gain is
modest and group-level---the individual-binary lift is small (AUC $+0.010$)---it is
co-informative with rather than superior to DSM-5, and it is carried by the biological
\emph{phenotype} (the immunometabolic archetype corner) rather than by any isolated
laboratory axis. ``Biology predicts functioning'' is therefore a phenotype-level
claim: the immunometabolic axis carries a credible adverse direction, but the
predictive object is the fuller archetype representation in which it is embedded. That
the eight-dimension map is nonetheless a \emph{sufficient} representation of the full raw
V0 indicator block ($143$ items, the representation benchmark's raw arm) for predicting
deterioration, and near-sufficient for recovery (the residual
is within-factor item-level compression, not a missing axis), shows the parsimony is
faithful rather than lossy. We report these modest and null findings as deliberately
as the positive ones; an atlas explicit about the gap between demonstrated scientific
validity and undemonstrated individual-level clinical utility is more useful than an
overstated one.

In sum, this atlas reframes a clinically actionable hypothesis with measured
coordinates: not ``which biotype is this patient?'' but ``where does this patient sit
on continuous, partly biological axes that diagnosis and severity do not capture, and
how durable is that position?'' The immunometabolic axis a patient keeps over time
forecasts the functional state they move toward---a group-level, transdiagnostic signal
that complements diagnosis and identifies the immunometabolic corner as a concrete
enrichment and monitoring target for prospective study. Separating biological burden from clinical
severity, and showing that the biological part is stable and prognostic, is the step
that makes mechanistic stratification in severe mental illness a measurable rather than
a rhetorical goal.

\subsection*{Limitations}

Three classes of limitation bound these claims. \emph{Internal validity only:} the map
is discovered and validated within FACE, with no external-cohort replication, and the
invariance partials that exist are documented rather than hidden (\Gfac{} re-anchors in
schizophrenia, which lacks the FAST functioning anchor; self-rated mania floors in
depression). Because FACE recruits at French expert centres, the cohort likely differs from
community, primary-care, first-episode and non-French samples in chronicity, comorbidity and
treatment intensity; the continuum geometry, the archetype simplex and the prognostic gradient
may shift in younger, first-episode or less-treated populations---a boundary external
replication must map.
\emph{Outcomes are scale trajectories, not recorded events:} prognostic endpoints are
state transitions defined from repeated functioning and severity scales, not registered
relapses or hospitalizations---a deliberate de-scope from event-level prognosis.
\emph{Measurement coverage and
horizon:} two axes are thin (mania has two indicators and is reported as data-limited;
substance use is a two-cohort axis), education is $40.8\%$ missing overall (and $74.5\%$
in schizophrenia)---a structured missingness the observed-likelihood design handles but
that limits that covariate---developmental risk scores as ``state'' partly because of
recall noise in childhood-adversity reports rather than true change, attrition is
informative (addressed by inverse-probability weighting), and the follow-up window is
two years across three visits. Establishing individual-level clinical utility
would require incident-event outcomes and external validation beyond this baseline
observational cohort.

\emph{Measurement precision is not uniform across axes.} A distinct limit is the
resolution the current instrument bank affords each axis. Six of the eight axes are
estimated with high reliability ($0.85$ and above), but two---mania/activation and
substance use---cannot exceed a reliability of about $0.45$ regardless of sample size
(maxima $0.408$ and $0.429$). This is a property of the instrument bank, not of missing
data: $95$--$96\%$ of patients have at least one item on these axes, but mania rests on
only two summary scores and substance use on four items, two of which are rare binary
flags that carry almost no information. Because the measurement model is missingness-native
and fully Bayesian, the consequence is graceful degradation rather than bias---an
under-instrumented axis returns a wide posterior rather than a false-precise point score,
so a barely-measured axis is never mistaken for a precisely located one. Two archetype
poles (\textbf{A0}, \textbf{A4}) are partly defined by the mania axis; these survive
nonetheless because an archetype coordinate is a group mean, and averaging over more than a
thousand patients per pole resolves the between-pole mania contrast to within $\approx0.02$
standard-deviation units even though any single patient's mania score is uncertain. The
model thus reports its own per-axis precision and names mania and substance use as the
priorities for the next data-collection wave---specifically graded severity instruments
(for example AUDIT- or DAST-style quantity--frequency measures) rather than additional
binary screens, which would add coverage without adding information. This sharpens the
external-validity boundary above: what a replication cohort should collect is not merely
more patients but better-instrumented mania and substance measures.
